\documentclass[11pt,a4paper]{article}
\usepackage[margin=1.7cm]{geometry}
\usepackage[T1]{fontenc}
\usepackage[utf8]{inputenc}
\usepackage[french]{babel}
\usepackage{lmodern,microtype,titlesec,enumitem,needspace}
\usepackage[hidelinks]{hyperref}
\input{glyphtounicode}
\pdfgentounicode=1
\renewcommand{\familydefault}{\sfdefault}
\pagestyle{empty}
\setcounter{secnumdepth}{0}
\setlength{\parindent}{0pt}
\setlength{\parskip}{3pt}
\titleformat{\section}{\large\bfseries}{}{0pt}{}
\titlespacing{\section}{0pt}{10pt}{5pt}
\setlist[itemize]{label=\textbullet,leftmargin=13pt,itemsep=2pt,topsep=3pt,parsep=0pt}
\newcommand{\role}[3]{\Needspace{5\baselineskip}\textbf{#1 | #2}\\#3\par}
\newcommand{\project}[2]{\Needspace{3\baselineskip}\textbf{#1}\\#2\par}
\hypersetup{pdftitle={Tamim Hmizi | Ingénieur logiciel | Cloud, DevOps et IA},pdfauthor={Tamim Hmizi},pdfsubject={Curriculum vitae professionnel},pdflang={fr}}
\begin{document}
{\fontsize{25}{28}\selectfont\textbf{Tamim Hmizi}}\par
{\large\textbf{Ingénieur logiciel | Cloud, DevOps et IA}}\par
Ariana | \href{mailto:tamimhmizi@icloud.com}{tamimhmizi@icloud.com} | +216 54 886 068\\
\href{https://www.linkedin.com/in/tamimhmizi/}{linkedin.com/in/tamimhmizi} | \href{https://github.com/tamim-hmizi}{github.com/tamim-hmizi}\\
\href{https://tamim-hmizi.github.io/Portfolio/}{tamim-hmizi.github.io/Portfolio/}

\section{Profil professionnel}
Ingénieur logiciel associant développement full stack, infrastructure cloud et intelligence artificielle (IA) appliquée. Fondateur d'Axynoxia et ingénieur full stack chez Bassetti Group. Je développe des applications web et mobiles, automatise leur déploiement et exploite leur infrastructure. Expérience sur Azure Stack Hub, AWS EKS, l'infrastructure as code, l'observabilité et l'intégration de grands modèles de langage (LLM).

\section{Compétences techniques}
\textbf{Développement :} TypeScript, JavaScript, Python, C\#, Java, PHP, Bash, SQL ; React, Next.js, Angular, Flutter ; Node.js, NestJS, Express, FastAPI, Flask, .NET, Spring Boot, Symfony.\\
\textbf{Cloud et DevOps :} Azure, Azure Stack Hub, AWS, EKS, AKS, OpenStack ; Docker, Kubernetes, Helm, Terraform, Ansible ; Jenkins, GitHub Actions, GitLab CI, CI/CD.\\
\textbf{IA et automatisation :} intégration de LLM, analyse de dépôts, recommandations de déploiement, automatisation Python, collecte web et intégration de workflows.\\
\textbf{Données et exploitation :} PostgreSQL, MongoDB, MySQL, Redis, Prisma, Firebase ; Prometheus, Grafana, SonarQube, Trivy, Nexus ; Linux, Windows Server, Commvault, Fortinet ; OAuth2, IAM, RBAC.

\section{Expérience professionnelle}
\role{Ingénieur full stack}{Bassetti Group, TEEXMA}{Avril 2026 à ce jour | Tunis}
\begin{itemize}
\item Développement d'applications de gestion de données techniques et scientifiques avec .NET, Angular et JavaScript sur la plateforme TEEXMA.
\item Maintenance de modules Delphi, conteneurisation avec Docker et contribution à l'automatisation des livraisons.
\end{itemize}

\role{Fondateur et ingénieur logiciel}{Axynoxia}{Janvier 2026 à ce jour | Tunis | \href{https://www.axynoxia.com}{www.axynoxia.com}}
\begin{itemize}
\item Pilotage des projets clients : cadrage, architecture, développement, CI/CD, déploiement et exploitation en production.
\item Réalisation de plateformes web, d'une application Flutter, de billetteries, d'un CMS immobilier et de boutiques Shopify ; intégration des paiements et des emails transactionnels.
\item Développement de produits internes ERP, billetterie, facturation, CRM et gestion énergétique, avec démonstrations hébergées et développements en cours.
\end{itemize}

\role{Consultant infrastructure}{RFC}{Octobre 2025 à mai 2026 | Ariana}
\begin{itemize}
\item Conception d'infrastructures hybrides Azure Stack Hub et Azure ; intervention sur les réseaux d'entreprise Fortinet et Cisco.
\item Administration des sauvegardes Commvault, validation des restaurations et analyse d'incidents ; automatisation des opérations Linux et Windows Server.
\item Développement d'exporteurs Python/Flask et de tableaux de bord Prometheus/Grafana pour la supervision des infrastructures et machines virtuelles.
\end{itemize}

\newpage
\section{Expériences précédentes}
\role{Ingénieur plateforme DevOps, PFE}{RFC}{Février 2025 à août 2025 | Ariana | \href{https://github.com/tamim-hmizi/QuickOps-AI}{Code public QuickOps}}
\begin{itemize}
\item Création de QuickOps, plateforme DevOps as a Service avec React, Express, FastAPI et MongoDB pour analyser, déployer et superviser des applications.
\item Intégration d'un moteur LLM analysant la structure des dépôts, les dépendances et les besoins d'exécution pour recommander un déploiement VM ou Kubernetes.
\item Automatisation des VM avec Terraform, Ansible et Docker Compose ; création de modèles de déploiement AKS pour les microservices.
\item Pipelines Jenkins avec SonarQube, Nexus, Prometheus et Grafana ; validation des deux parcours sur deux applications e-commerce MERN.
\end{itemize}

\role{Stagiaire cloud et full stack}{RFC}{Juillet 2024 à août 2024}
\begin{itemize}
\item Développement d'une application de gestion des stagiaires et déploiement de services conteneurisés sur AWS EKS.
\item CI/CD avec Jenkins, Docker, Terraform et GitHub Actions ; contrôles IAM et supervision Prometheus/Grafana.
\end{itemize}

\role{Stagiaire développeur web}{ESPRIT}{Juillet 2023 à août 2023}
\begin{itemize}
\item Plateforme de recherche de cours : collecte Python sur des sites éducatifs officiels, normalisation des métadonnées, filtres et recommandations.
\end{itemize}

\role{Stagiaire ingénieur réseau, C-NOC}{Tunisie Telecom}{Août 2021 à septembre 2021}
\begin{itemize}
\item Supervision des liaisons d'entreprise au centre d'exploitation réseau B2B ; assistance au diagnostic d'incidents de routage et de disponibilité.
\end{itemize}

\section{Projets cloud, DevOps et logiciels}
\project{Supervision Azure Stack Hub | RFC}{Exporteurs Python/Flask exploitant les API Admin avec OAuth2 pour collecter la télémétrie infrastructure et VM ; Docker, Prometheus et Grafana pour les alertes et tableaux de bord.}
\project{DevOps CI/CD Hub | ESPRIT}{Pipelines Jenkins pour Angular et Spring Boot, contrôles qualité SonarQube, analyse de sécurité Trivy et gestion des artefacts Nexus.}
\project{Laboratoire cloud privé OpenStack | ESPRIT}{Cloud multi-nœuds : réseaux virtuels, sous-réseaux, calcul, stockage et connectivité VPN pour l'hébergement d'applications.}
\project{E-commerce MERN | Projet personnel}{Boutique React/Redux, API Express et MongoDB, authentification JWT, hachage Argon2, administration des commandes et gestion des médias.}
\project{New Journey et Jektis Travel | Projets personnels}{Plateforme de conseil Symfony 7/PHP 8.2 avec Doctrine, Turbo, Stimulus et RBAC ; plateforme de voyage PHP avec réservations, collecte de prix et notifications email.}

\section{Formation et langues}
\textbf{Diplôme d'ingénieur en informatique | ESPRIT} | 2020 à 2025 | Ariana\\
Formation générale en génie logiciel ; spécialisation Architecture Cloud et DevOps (ArcTIC).\\
\textbf{Langues :} arabe (langue maternelle), anglais (B2), français (B2).

\newpage
\section{Projets clients}
\project{B'moovd Ticketing | \href{https://bmoovd-shop.com}{bmoovd-shop.com}}{Billetterie Next.js et tRPC avec MongoDB/Prisma, paiements intégrés, billets QR signés, contrôle d'accès par caméra et envoi de PDF par email.}
\project{Application B'moovd | \href{https://apps.apple.com/de/app/bmoovd/id6736364332}{App Store} | \href{https://play.google.com/store/apps/details?id=com.bmoovd.app}{Google Play}}{Application Flutter de découverte de lieux et de réservation : Firebase Auth, Firestore, Cloud Functions, stockage et intégration des paiements.}
\project{Roehrdanz Immobilien | \href{https://rimmobilien-test.de/}{rimmobilien-test.de}}{CMS immobilier React, Express et MongoDB avec administration des annonces, gestion des images, droits d'accès et recherche de biens.}
\project{Confortex et Vie Sublime | \href{https://confortex.shop/}{confortex.shop} | \href{https://www.viesublime.tn/}{viesublime.tn}}{Boutiques Shopify avec thèmes personnalisés, catalogues par collections, contenu éditorial, inscription newsletter et informations sur les showrooms.}
\project{Habiter Mieux | \href{https://habiter-mieux.info/}{habiter-mieux.info}}{Site d'installation solaire avec présentation des services et parcours de demande de devis relié au traitement des demandes client.}
\project{OMHY Entertainment | \href{https://omhyfamily.com/}{omhyfamily.com}}{Site média et portfolio React, interfaces Chakra UI, animations GSAP, services de données MongoDB/MySQL et pipelines de livraison Jenkins.}
\project{NexGen Studios}{Site vitrine Next.js App Router, conteneurisé avec Docker et servi par Nginx.}

\section{Démonstrations de produits Axynoxia}
\project{Axynoxia ERP | \href{https://erp.products.axynoxia.com}{erp.products.axynoxia.com}}{ERP modulaire NestJS, Prisma, PostgreSQL et Redis : découverte des modules, gestion des accès et déploiements dédiés par client. Développement en cours.}
\project{Tickety | \href{https://tickets.products.axynoxia.com}{tickets.products.axynoxia.com}}{Billetterie React/Vite et TypeScript en monorepo pnpm, avec expiration des commandes, traitement des remboursements et tests de bout en bout.}
\project{Axynoxia Energie | \href{https://energie.products.axynoxia.com}{energie.products.axynoxia.com}}{Démonstration de gestion énergétique solaire et stockage batterie en TypeScript : simulateur de centrale, télémétrie, alarmes et tableaux de supervision.}
\project{ProInvoice | \href{https://invoice.products.axynoxia.com}{invoice.products.axynoxia.com}}{Application de facturation React et Express/TypeScript sur MongoDB : double authentification, exports PDF/Excel, emails et documentation API Swagger.}
\project{Axynoxia CRM et Leasing | \href{https://crm-leasing.products.axynoxia.com}{crm-leasing.products.axynoxia.com}}{Démonstration MERN de gestion des contrats de leasing, calculs financiers, droits par rôle, interface arabe RTL et traitements asynchrones.}

\vfill
\textbf{Code, projets et contact :} \href{https://tamim-hmizi.github.io/Portfolio/}{tamim-hmizi.github.io/Portfolio/}
\end{document}
